\begin{table}
\centering
\caption{Lista de valores de pKa de algunos aminoácidos en base a su pH}
\label{tab:pka_aminoácidos}
\begin{threeparttable}
{\setstretch{2}
\begin{tabular}{cccc}\hline

    \textbf{Aminoácido} & \textbf{$pK_{a}$ de grupo carboxilo $COOH$} & \textbf{$pK_{a}$ de grupo amino $NH_{3}^{+}$} & $pK_{a}$ de cadena lateral\tnote{1} \\ \hline
    
    Ácido aspártico & 2.09 & 9.82 & 3.86 \\
    Ácido glutámico & 2.19 & 9.67 & 4.25 \\
    Alanina & 2.34 & 9.69 & - \\
    Arginina & 2.17 & 9.04 & 12.48 \\
    Asparagina & 2.02 & 8.84 & - \\
    Cisteína & 1.92 & 10.46 & 8.35 \\
    Fenilalanina & 2.16 & 9.18 & - \\
    Glicina & 2.34 & 9.60 & - \\
    Glutamina & 2.17 & 9.13 & - \\
    Histidina & 1.82 & 9.17 & 6.04 \\
    Isoleucina & 2.36 & 9.68 & - \\
    Leucina & 2.36 & 9.60 & - \\
    Lisina & 2.18 & 8.95 & 10.79 \\
    Metionina & 2.28 & 9.21 & - \\
    Prolina & 1.99 & 10.60 & - \\
    Serina & 2.21 & 9.15 & - \\
    Tirosina & 2.20 & 9.11 & 10.07 \\
    Treonina & 2.63 & 9.10 & - \\
    Trptófano & 2.38 & 9.39 & - \\
    Valina & 2.32 & 9.62 & - \\
    
    \hline
\end{tabular}
}
\begin{tablenotes} % ex: \tntote{1} ... \item[1]
    \item[1] Los valores que no aparecen (-) ...
\end{tablenotes}
\end{threeparttable}
\end{table}